\section{Attribution, editions and the separate questions}
The proved source conclusions covered here are Propositions 2.2--2.4, Theorem 3.1, Remark 3.2, Corollaries 3.3--3.4, Example 3.5, Proposition 3.6, Corollary 3.7, Example 3.8, the uniqueness implication at the beginning of Section 4 and all claims in its explicit finite example. The exact source is Gabe Cunningham, \emph{Chiral polytopes whose smallest regular cover is a polytope}, the complete eight-page author manuscript headed 25 July 2023. Publication metadata is Journal of Combinatorial Theory, Series A 204 (2024), 105839. The full publisher PDF was not acquired, so this is not a certification of identical wording in the version of record. The earlier publisher metadata notes a correction concerning Corollary 3.3; our proof binds the precise embedded-subgroup statement in the checked author edition.

The common group calculations are elementary reformulations of arguments already used in this literature and the accepted proof bank. The source's own propositions are not assumed to prove its main theorem. The string $C^+$ correspondence and depth-two section theorem remain retained geometric inputs; the square-tessellation model and the specified atomic-family classification remain additional retained inputs. The finite groups are established by the complete required GAP program. The stated Lean certificates verify their exact generic interfaces, not these source-specific geometric or computational instantiations.

The source also asks whether nonpolytopality of the smallest regular cover always forces multiple minimal regular polytope covers, and conjectures the analogous phenomenon for all five tight families in Example 3.5. These are not proved by the arguments of this application. Later admitted project results concerning them have their own proofs, reviews and quantifiers: the general statement allowing arbitrary regular covers, finite rank-four results, and interval-extension results must not be conflated with a theorem about finite covers in every rank. Their status is recorded in the accompanying question map. Those later proofs are not used as unpriced premises here, and no new question credit is assigned to this canonical edition.

The statement requiring finite covers in every rank remains unresolved in the closed corpus. Results for arbitrary independent markings do not automatically apply to the monodromy groups of actual polytopes. This application preserves the source's complete proved conclusions without importing a stronger claim from that separate investigation.
